\begin{Answer}{3.5}
The prime numbers less than $10$ are $2$, $3$, $5$, and $7$.
But the problem asked for the {\em set of} prime numbers less than $10$.
Therefore, the answer is $\{2, 3, 5, 7\}$.
If you were asked to {\em list} the prime numbers less than $10$,
an appropriate answer would have been $2,3,5,7$ (but that is not what was asked).
The cardinality of the set is 4. That is, $|\{2, 3, 5, 7\}|=4$.
\end{Answer}
\begin{Answer}{3.8}
6; 5; 6; $A$ and $C$ represent the same set. That is, $A=C$.
\end{Answer}
\begin{Answer}{3.11}
$\infty$; $\infty$.  You might think it is $\infty/2$, but you can't do arithmetic with $\infty$ since it
isn't a number. Without getting too technical, although $\BBZ^+$ seems to have about half as many elements
as $\BBZ$, it actually doesn't.  It has the exact same number: $\infty$. ; $0$.
\end{Answer}
\begin{Answer}{3.14}
$\{2a : a\in\integers \}$ and $\{\ldots, -4, -2, 0, 2, 4, \ldots\}$.
\end{Answer}
\begin{Answer}{3.17}
$\BBQ = \{a/b : a,b\in\integers,b\neq 0 \}$.
\end{Answer}
\begin{Answer}{3.20}
(a) Yes.
(b) Yes.  $A$ is a proper subset since $25$, for instance, is in $S$ but not in $A$.
(c) Yes.  Every set is a subset of itself.
(d) No.  No subset is a {\em proper subset} of itself.
(e) No.  $25\in S$, but $25\not\in A$.
\end{Answer}
\begin{Answer}{3.21}
(a) yes.  Any number that is divisible by 6 is divisible by 2.;
(b) yes. Any number that is divisible by 6 is divisible by 3.;
(c) no.  $4\in B$, but $4\not\in A$.;
(d) no.  $4\in B$, but $4\not\in C$.;
(e) no.  $3\in C$, but $3\not\in A$.;
(f) no.  $3\in C$, but $3\not\in B$.
\end{Answer}
\begin{Answer}{3.24}
 We will use the result of example
\ref{exa:subsets3element}.
A subset of $\{a, b, c, d\}$
either contains $d$ or it does not.
Since the subsets of $\{a, b, c\}$ do not contain $d$, we simply list all the subsets of $\{a, b, c\}$ and then to each one of them  we add $d$. This gives

$\begin{array}{lllllll} S_1 & = & \varnothing & \hspace{1cm} &  S_9 & = & \{d\} \\
S_2 & = & \{a\} & \hspace{1cm} &  S_{10} & = & \{a, d\} \\
S_3 & = & \{b\}& \hspace{1cm} &  S_{11} & = & \{b, d\} \\
S_4 & = & \{c\} & \hspace{1cm} &  S_{12} & = & \{c, d\} \\
S_5 & = & \{a, b\}& \hspace{1cm} &  S_{13} & = & \{a, b, d\} \\
S_6 & = & \{b, c\}& \hspace{1cm} &  S_{14} & = & \{b, c, d\}  \\
S_7 & = & \{a, c\} & \hspace{1cm} &  S_{15} & = & \{a, c, d\} \\
S_8 & = & \{a, b, c\} & \hspace{1cm} &  S_{16} & = & \{a, b, c, d\}  \\
\end{array}   $
\end{Answer}
\begin{Answer}{3.27}
Based on the answer to Exercise~\ref{exa:subsets4element}, we have that
$P(\{a, b, c, d\})=\{ \varnothing,
 \{a\}, \{b\},  \allowbreak\{c\}, \allowbreak\{a, b\},\allowbreak\{b, c\},
\allowbreak\{a, c\}, \allowbreak\{a, b, c\},\allowbreak  \{d\},\allowbreak \{a, d\},
 \allowbreak\{b, d\},\allowbreak \{c, d\}, \allowbreak \{a, b, d\},\allowbreak \{b, c, d\},
\allowbreak\{a, c, d\}, \allowbreak \{a, b, c, d\}
\}$.
Notice that a list of these 16 sets not separated by commas and not enclosed in $\{\}$ is
not correct.  It may have the correct {\em content}, but it is not in the proper {\em form}.
\end{Answer}
\begin{Answer}{3.29}
(a) By Theorem~\ref{thm:powerSetSize}, $|P(A)|=2^4=16$.
(b) Similarly, $|P(P(A))|=2^{16}=65536$.
(c) This is just getting a bit ridiculous, but the answer is $|P(P(P(A)))|=2^{65536}$.
\end{Answer}
\begin{Answer}{3.30}
Applying Theorem~\ref{thm:powerSetSize}, it is not too hard to see that the power set will be
twice as big after a single element is added.
\end{Answer}
\begin{Answer}{3.33}
$\BBZ$, or the set of (all) integers.
\end{Answer}
\begin{Answer}{3.36}
$\varnothing$.
\end{Answer}
\begin{Answer}{3.39}
$A$; $B$.
\end{Answer}
\begin{Answer}{3.43}
$B$; $A$.
\end{Answer}
\begin{Answer}{3.47}
Since no integer is both even and odd, $A$ and $B$ are disjoint.
\end{Answer}
\begin{Answer}{3.49}
Let $E$ be the set
of all English speakers, $S$ the set of Spanish speakers and $F$ the set of

\noindent
\begin{tabular}{p{.55\textwidth}p{.38\textwidth}}
French speakers in our group.
We fill-up the  Venn diagram (to the right) successively. In the
intersection of all three we put 3. In the region common to $E$ and
$S$ which is not filled up we put $5 - 3 = 2$. In the region common
to $E$ and $F$ which is not already filled up we put $5 - 3 = 2$. In
the region common to $S$ and $F$ which is not already filled up, we
put $7 - 3 = 4.$ In the remaining part of $E$ we put $8 - 2 - 3 - 2
= 1,$ in the remaining part of $S$ we put $12 - 4 - 3 - 2 = 3$, and
in the remaining part of $F$ we put $10 - 2 - 3 - 4 = 1$. Therefore,
 $1 + 2 + 3 + 4 + 1 + 2 + 3 = 16$ people speak at least one of these languages.
&
\begin{pspicture}(-3, 3)(3,2.5)
\psset{unit=1.5pc}
\pscustom[fillstyle=solid,fillcolor=white]{\psline(-4.25,4.25)(-4.25,-2.5)(4.25,-2.5)(4.25,4.25)(-4.25,4.25)}
\pscircle(-1,0){2}
\pscircle(1,0){2}
\pscircle(0,1.41421356){2}
\rput(2,0){\small 3}
\rput(-2, 0){\small 1}
\rput(0,0.707106781){\small 3}
\rput(0, 2.3){\small 1}
\rput(0,-0.9){\small 2}
\rput(-1.43,1.43){\small 2}
\rput(1.43,1.43){\small 4}
\rput(-3.8, 0){\small E}
\rput(3.8,0){\small S}
\rput(0, 3.8){\small F}
\end{pspicture}
\end{tabular}
\end{Answer}
\begin{Answer}{3.52}
$A\times B = \{(1,3), (2,3), (3,3), (4,3)\}$.
\end{Answer}
\begin{Answer}{3.55}
$A^2=\{(0,0), (0,1), (1,0), (1,1)\}$.

\noindent
$A^3=\{(0,0,0), (0,1,0), (1,0,0), (1,1,0), (0,0,1), (0,1,1), (1,0,1), (1,1,1)\}$.
\end{Answer}
\begin{Answer}{3.58}
(a) $10*50=500$
(b) $10*20=200$
%(c) $10*10=100$
(c) $50*50*50=125,000$
(d) $10*50*20=10,000$
\end{Answer}
\begin{Answer}{3.59}
\begin{solevalenum}
\item
Although it is on the right track, this solution has several problems.  First, it would be better to make
it more clear that the assumption is that {\em both} $A$ and $B$ are not empty.
But the bigger problem is the statement `$(a,b)$ is in the cross product'.
The problem is that $a$ and $b$ are not defined anywhere.
Saying `where $a\in A$ and $b\in B$' earlier does not guarantee that there is such an $a$ or $b$.
The proof needs to say something along the lines of
`Since $A$ and $B$ are not empty, then there exist some $a\in A$ and $b\in B$.  Therefore
$(a,b)\in A\times B$\ldots'
\item
This one is way off.  The proof is essentially saying
`Notice that $p\rightarrow q$. Therefore $q\rightarrow p$.'  But these are not equivalent
statements.  Although it is true that if both $A$ and $B$ are the empty set, then $A\times B$ is
also the empty set, this does not prove that both $A$ and $B$ must be empty in order for $A\times B$
to be empty.  In fact, this isn't the correct conclusion.
\item
The conclusion is incorrect, as is the proof.
The problem is that the negation of
`both $A$ and $B$ are empty' is `it is not the case that both $A$ and $B$ are empty'
or `at least one of $A$ or $B$ is not empty,' which is not the same thing as
`neither $A$ nor $B$  is empty.'  So although the proof seems to be correct, it is not.
The reason it seems almost correct is that except for this error, the rest of the proof follows
proper proof techniques.  Unfortunately, all it takes is one error to make a proof invalid.
\item
This is a correct conclusion and proof.
\end{solevalenum}
\end{Answer}
\begin{Answer}{3.63}
\begin{pfevalenum}
\item
This solution has several problems.
\begin{enumerate}
\item
$x\in \{A-B\}$ means `$x$ is an element of the set containing $A-B$,' not
`$x$ is an element of $A-B$.'  What they meant was `$x\in A-B$.'
\item
At the end of the first sentence, `$x$ is not $\in B$' mixes mathematical
notation and English in a strange way.
This should be either `$x\not\in B$' or `$x$ is not in $B$.'
\item
In the second sentence,
the phrase  `$x\in A$ and $\overline{B}$' is a strange mixture of math and English
that is potentially ambiguous.  It should be rephrased as something like
`$x\in A$ and $x\in\overline{B}$' or `$x$ is in both $A$ and $\overline{B}$.'
\item
Finally, what has been shown here is that $A-B\subseteq A\cap\overline{B}$.
This is only half of the proof.  They still need to prove that
$A\cap\overline{B}\subseteq A-B$.
\end{enumerate}
\item
Overall, this proof is very confusing and unclear.  More specifically,
\begin{enumerate}
\item
This is an attempt at working through what each set is by using the definitions.
That would be fine except for two things.
First, they were asked to give a set containment proof.
Second, the wording of the proof is confusing and hard to follow.
I do not come away from this with a sense that anything has been proven.
\item
They are not using the terminology properly. The terms `universe' or `universal set' would be appropriate, but not `universal' on its own (used twice).  Similarly, what does the phrase
`all intersection part' mean?
Also, a set doesn't `return' anything.  A set is just a set. It contains elements, but
it doesn't `do' anything.
\end{enumerate}
\item
This proof contains a lot of correct elements.  In fact, the first half is on the right track.
However, they jumped from $x\in A$ and $x\not\in B$ to $x\in A\cap\overline{B}$.  Between
these statements they should say something like
`$x\not\in B$ is equivalent to $x\in\overline{B}$' since the latter statement is really needed
before they can conclude that $x\in A\cap\overline{B}$.
Also, it would be better if they had `by the definition of intersection'
before or after the statement $x\in A\cap\overline{B}$.
Finally, it would help clarify the proof if the end was something like
'We have shown that whenever $x\in A-B$, $x\in A\cap\overline{B}$.  Thus, $A-B\subseteq A\cap\overline{B}$.'

The second half of the proof starts out well, but has serious flaws.
The statement 'This means that $x\in A$ and $x\not\in B$'
should be justified by the definitions of complement and intersection, and might even involve two steps.
This is the same problem they had in the first half of the proof.
More serious is the statement `which is what we just proved in the previous statement'.  What exactly does
that mean?  It is unclear how `what we just proved' immediately leads us to the conclusion that
$A-B = A\cap\overline{B}$.  First we need to establish that $x\in A-B$ based on the previous statements (easy).
Then we can say that $A\cap\overline{B}\subseteq A-B$.  Finally, we can combine this with the first part
of the proof to say that $A-B = A\cap\overline{B}$.

In summary, the first half is pretty good.  It should at least make the connection between
$x\not\in B$ and $x\in\overline{B}$.  The other suggestions clarify the proof a little,
but the proof would be O.K. if they were omitted.  The second half is another story.  It doesn't really
prove anything, but instead makes a vague appeal to something that was proven before.  Not only is
what they are referring to unclear, but how the proof of one direction is related to the proof of
the other direction is also unclear.
\end{pfevalenum}
\end{Answer}
\begin{Answer}{3.65}
$x\in C$;
$x\in B$;
definition of union;
$(x\in B \wedge x\in C)$;
distributive law (the logical one);
$(x\in A\cap C) $;
definition of intersection;
definition of union.
\end{Answer}
\begin{Answer}{3.78}
(a) 45; (b) 8; (c) 3; (d) 6; (e) 0; (f) 7; (g) 7; (h) 7; (i) 11.
\end{Answer}
\begin{Answer}{3.81}
Since every integer is either even (of the form $2k$) or odd (of
the form $2k + 1$) we have two possibilities:
$$\begin{array}{lllll}(2k)^2 & = & 4k^2&\equiv&0\pmod 4, \mbox{or}\\
(2k + 1)^2 & = & 4(k^2 + k) + 1&\equiv&1\pmod 4.
  \end{array}$$
Thus, $n^2$ has remainder $0$ or $1$ when divided by $4$.
\end{Answer}
\begin{Answer}{3.86}
Observe that $3^{2n + 1} \equiv 3\cdot 9^n \equiv 3\cdot
2^n$ mod 7 and $2^{n + 2} \equiv 4\cdot 2^n$ mod 7. Hence
$ 3^{2n + 1} + 2^{n + 2} \equiv 3\cdot 2^n+4\cdot 2^n\equiv 7\cdot 2^n \equiv 0 \  \bmod  \ 7,$
for all
natural numbers $n.$
\end{Answer}
\begin{Answer}{3.89}
 $2^1 \equiv 2, 2^2
\equiv 4, 2^3 \equiv 1$ mod 7, and this cycle of three repeats.
Thus $2^k - 5$ can leave only remainders
$4=(2-5)\bmod 7$, $6=(4-5)\bmod 7$, or $3=(1-5\bmod 7)$ upon
division by $7$, none of which are 1.
\end{Answer}
\begin{Answer}{3.93}
-15; -7; 9; 13; 21.  Notice that it is every 4th number along the number line, both in the
positive and negative directions.
\end{Answer}
\begin{Answer}{3.99}
$\gcd(524,118)=2$ as demonstrated here:

\begin{tabular}{|l|l|l|l|}
\hline
step&a&b&r\\
\hline
\hline
1&524&118&52\\
\hline
2&118&52&14\\
\hline
3&52&14&10\\
\hline
4&14&10&4\\
\hline
5&10&4&2\\
\hline
6&4&2&0\\
\hline
7&2&0&0\\
\hline
\end{tabular}
\end{Answer}
\begin{Answer}{3.102}
You should have computed that $\gcd(867,5309)=1$ and therefore
concluded that they are relatively prime.
\end{Answer}
\begin{Answer}{3.105}
(1) 9; (2) 10; (3) 9; (4) 10; (5) 9; (6) 9.
\end{Answer}
\begin{Answer}{3.114}
$f(x)=x\bmod 2$ works.  The domain is $\BBZ$, and the codomain can be a variety
of things.  $\BBZ$, $\naturals$, and $\{0,1\}$ are the most obvious choices.
Note that we can pick any of these since the only requirement of the codomain is
that the range is a subset of it.  On the other hand, $\BBR$, $\BBC$ and $\BBQ$
could also all be given as the codomain, but they wouldn't make nearly as much sense.
\end{Answer}
\begin{Answer}{3.118}
(a) F.  Consider $f(x)=\lfloor x \rfloor$ from $\reals$ to $\integers$.
(b) F.  Consider $f(x)=x^2$ from $\reals$ to $\reals$ which is not one-to-one.
(c) T.  See Theorem~\ref{thm:size_domain_image_injections_surjections}.
(d) F.  $f$ maps $1$ to two different values, so it isn't a function.
(e) T.  We previously showed it was onto, and it isn't difficult to see that it is one-to-one.
(f) F. $f$ is not onto, but it is one-to-one.
(g) T.  By definition of range, it is a subset of the codomain.
(h) F.  We have seen several counter examples to this.
(i) F.  If $a=2$ and $b=0$, the odd numbers are not in the range.
(j) F.  Same counterexample as the previous question.
(k) T.  The proof is similar to several previous proofs.
\end{Answer}
\begin{Answer}{3.124}
Let $y=7x+2$.  Then $7x=y-2$, so $x=(y-2)/7$.  Thus, $f^{-1}(x)=(x-2)/7$
(or $\frac{x}{7}-\frac{2}{7}$).
\end{Answer}
\begin{Answer}{3.127}
$(f\circ g)(x)=f(x/2)=\lfloor x/2\rfloor$, and
$(g\circ f)(x)=g(\lfloor x\rfloor) = (\lfloor x\rfloor)/2$.
\end{Answer}
\begin{Answer}{3.131}
(a) F.  $f$ might not be onto--e.g. if $a=2$ and $b=0$.
(b) F. Same reason as the previous question.
(c) T. Since over the reals, $f$ is one-to-one and onto.
(d) F. There are several problems.  First, $x^2$ may not even have an inverse depending on
the domain (which was not specified).  Second, even if it had an inverse, it certainly wouldn't
be $1/x^2$.  That's its reciprocal, not its inverse.  Its inverse would be $\sqrt{x}$ (again,
assuming the domain was chosen so that it is invertible).
(e) F.  This is only true if $n$ is odd.
(f) F.  $\sqrt{2}\not\in\naturals$, so not only is it not invertible, it can't even be defined on $\naturals$.
(g) T.  The $n$th root of a positive number is defined for all positive real numbers, so the function is
well defined.  It is not too difficult to convince yourself that the function is both one-to-one and onto
when restricted to positive numbers, so it is invertible.
(h) T.  In both cases you get $1/x^2$.
(i) F.  $(f\circ g)(x)=f(x+1)=(x+1+1)^2=(x+2)^2=x^2+4x+4$, and $(g\circ f)(x)=g((x+1)^2)=(x+1)^2+1=x^2+2x+2$,
which are clearly not the same.
(j) F. $(f\circ g)(x)=\lceil x\rceil$, and $(g\circ f)(x)=\lfloor x\rfloor$.
(We'll leave it to you to see why this is the case.)
(k) F.  Certainly not.  $f(3.5)=3$, but $g(3)=3$, not $3.5$.
(l) T. With the restricted domain, they are indeed inverses.
\end{Answer}
\begin{Answer}{3.134}
We never said it was {\em always} wrong to work both sides of an equation.  If you are working
on an equation that you know to be true, there is absolutely nothing wrong with it.
It is a problem only when you are starting with something you don't know to be true.
In this case, we know that $2a-3=2b-3$ is true given the assumption made.  Therefore, we are
free to `work both sides'.
\end{Answer}
\begin{Answer}{3.135}
Let $a,b\in \BBR$.  If $f(a)=f(b)$, then $5a=5b$.  Dividing both sides by $5$, we get $a=b$.
Thus, $f$ is one-to-one.
\end{Answer}
\begin{Answer}{3.138}
Notice that $f(4.5)=f(4)=4$, so clearly $f$ is not one-to-one.
(Your proof may involve different numbers, but should be this simple.)
\end{Answer}
\begin{Answer}{3.141}
Notice that if $y=2x+1$, then $y-1=2x$ and $x=(y-1)/2$.
Let $b\in\reals$.  Then $f((b-1)/2)=2((b-1)/2)+1 = b-1+1=b$.
Thus, every $b\in\reals$ is mapped to by $f$, so $f$ is onto.
\end{Answer}
\begin{Answer}{3.144}
Since the floor of any number is an integer, there is no $a$ such that $f(a)=4.5$ (for instance).
Thus, $f$ is not onto.
\end{Answer}
\begin{Answer}{3.145}
(a) {\bf $f$ is not one-to-one.}  See Example~\ref{exa:xsquarenotoneone} for a proof.
(b) The same proof from Example~\ref{exa:xsquarenotoneone} works over the reals.
But I guess it doesn't hurt to repeat it:  Since $f(-1)=f(1)=1$, {\bf $f$ is not one-to-one.}
(c) Let $a,b\in\naturals$.  If $f(a)=f(b)$,
that means $a^2=b^2$. Taking the square root of both sides, we obtain $\sqrt{a^2}=\sqrt{b^2}$, or
$|a|=|b|$  (if you didn't remember that
$\sqrt{x^2}=|x|$, you do now).
But since $a,b\in\naturals$, $|a|=a$ and $|b|=b$.  Thus, $a=b$.  Thus, {\bf $f$ is one-to-one.}
\end{Answer}
\begin{Answer}{3.146}
If $f(a)=f(b)$, $3a-5=3b-5$.  Subtracting 5 from both sides and then dividing both sides by $3$,
we get $a=b$.  Thus, $f$ is one-to-one.
If $b\in\reals$, notice that $f((b+5)/3)=3((b+5)/3)-5 = b+5-5=b$, so there is some value that maps
to $b$.  Therefore, $f$ is onto.  Since $f$ is one-to-one and onto, it has an inverse.
To find the inverse, we
let $y=3x-5$.  Then $3x=y+5$, so $x=(y+5)/3$.  Thus, $f^{-1}(x)=(x+5)/3$
(or $\frac{x}{3}+\frac{5}{3}$).
\end{Answer}
\begin{Answer}{3.147}
(a)
Notice that if $f(a)=f(b)$, then $a-7=b-7$ so $a=b$.  Thus, $f$ is one-to-one.
Also notice that for any $b\in\integers$, $f(b+7)=b+7-7=b$, so $f$ is onto.
(b) Since $g(1)=g(-1)=1$, $g$ is not one-to-one.  Also notice that there is no integer $a$ such that
$g(a)=a^4=5$, so $g$ is not onto.
(c) If $h(a)=h(b)$, then $3a=3b$ so $a=b$.  Thus, $h$ is one-to-one.
But there is no integer $a$ such that $h(a)=3a=1$, so $a$ is not onto.
(d) Notice that $r(0)=\lfloor 0/2\rfloor=\lfloor 0\rfloor = 0$ and
$r(1)=\lfloor 1/2\rfloor=\lfloor 0\rfloor = 0$, so $r$ is not one-to-one.
But for any integer $b$, $r(2b)=\lfloor 2b/2\rfloor=\lfloor b\rfloor = b$, so $r$ is onto.
\end{Answer}
\begin{Answer}{3.155}
The following three cases probably make the most sense: When $a=b$, when $a<b$ and when $a>b$.
These make sense because these are likely different cases in the code.
Mathematically, we can think of it as follows.  The possible inputs are from the set $\BBZ\times\BBZ$.
The partition we have in mind is $A=\{(a,a) : a\in \BBZ\}$, $B=\{(a,b) : a,b\in \BBZ, a<b\}$, and
$C=\{(a,b) : a,b\in \BBZ, a>b\}$.  Convince yourself that these sets form a partition of $\BBZ\times\BBZ$.
That is, they are all disjoint from each other and $\BBZ\times\BBZ =A\cup B\cup C$.

Alternatively, you might have thought in terms of $a$ and/or $b$ being positive, negative, or $0$.
Although that may make some sense, given that we are comparing $a$ and $b$ with each other, it probably
doesn't matter exactly what values $a$ and $b$ have (i.e. whether they are positive, negative, or $0$),
but what values they have {\em relative to each other}.  That is why the first answer is much better.
With that being said, it wouldn't hurt to include several tests for each of our three cases that involve
various combinations of positive, negative, and zero values.
\end{Answer}
\begin{Answer}{3.156}
Did you define two or more subsets of $\BBZ$?
Are they all non-empty?
Do none of them intersect with each other?
If you take the union of all of them, do you get $\BBZ$?
If so, your answer is correct!  If not, try again.
\end{Answer}
\begin{Answer}{3.158}
Since $\BBR = \BBQ \cup \BBI$ and $\BBQ \cap \BBI = \varnothing$, $\{\BBQ,\BBI\}$ is a partition of $\BBR$.
Hopefully this comes as no surprise.
\end{Answer}
\begin{Answer}{3.163}
$R$ is a subset of $\BBZ\times\BBZ$, so it is a relation.
By the way, this relation should look familiar.
Did you read the solution to Exercise~\ref{exer:maxPart}?
\end{Answer}
\begin{Answer}{3.164}
Is it a subset of $\BBZ^+\times\BBZ^+$?  It is.  So it is a relation on $\BBZ^+$.
\end{Answer}
\begin{Answer}{3.166}
(a) $T$ is {\bf not} reflexive since you cannot be taller than yourself.
(b) $N$ is reflexive because everybody's name starts with the same letter as their name does.
(c) $C$ is reflexive because everybody have been to the same city as they have been in.
(d) $K$ is {\bf not} reflexive because you know who you are,
so it is not the case that you don't know who you are.  That is, $(a,a)\not\in K$ for any $a$.
(e) $R$ is {\bf not} reflexive because $(\mbox{Donald Knuth}, \mbox{Donald Knuth})$
(for instance) is not in the relation.
\end{Answer}
\begin{Answer}{3.168}
(a) $T$ is {\bf not} symmetric since if $a$ is taller than $b$, $b$ is clearly not taller than $a$.
(b) $N$ is symmetric since if $a$'s name starts with the same letter as $b$'s name, clearly $b$'s name
starts with the same letter as $a$'s name.
(c) $C$ is symmetric since it is worded such that it doesn't distinguish between the first and second
item in the pair.  In other words, if $a$ and $b$ have been to the same city, then $b$ and $a$ have
been to the same city.
(d) $K$ is {\bf not} symmetric since $(\mbox{David Letterman},\mbox{Chuck Cusack})\in K$, but
$(\mbox{Chuck Cusack},\mbox{David Letterman})\not\in K$.
(e) $R$ is not symmetric since $(\mbox{Barack Obama}, \mbox{George W. Bush})\in R$, but
(George W. Bush, Barack Obama)$\not\in R$.
\end{Answer}
\begin{Answer}{3.170}
(a) Just knowing that $(1,1)\in R$ is not enough to tell either way.
(b) On the other hand, if $(1,2)$ and $(2,1)$ are both in $R$, it is clearly
{\bf not} anti-symmetric.
\end{Answer}
\begin{Answer}{3.171}
This is just the contrapositive of the original definition.
\end{Answer}
\begin{Answer}{3.172}
(a) $T$ is anti-symmetric since whenever $a\neq b$, if $a$ is taller than $b$, then $b$ is not taller than $a$,
so if $(a,b)\in T$, then $(b,a)\not\in T$.
(b) $N$ is {\bf not} anti-symmetric since $(\mbox{Bono},\mbox{Boy George})$ and $(\mbox{Boy George},\mbox{Bono})$
are both in $N$.
(c) $C$ is {\bf not} anti-symmetric since $(\mbox{Bono},\mbox{The Edge})$ and $(\mbox{The Edge},\mbox{Bono})$
are both in $C$ (since they have played many concerts together, they have certainly been in the same
city at least once).
(d) Since
both $(\mbox{Dirk Benedict},\mbox{Jon Blake Cusack 2.0})$ and
$(\mbox{Jon Blake Cusack 2.0},\mbox{Dirk Benedict})$ are in $K$,
$K$ is {\bf not} anti-symmetric.
%(d) $K$ is {\bf not} anti-symmetric because both $(\mbox{Brian Yurk},\mbox{Chuck Cusack})$ and
%$(\mbox{Chuck Cusack},\mbox{Brian Yurk})$ are in $K$.
(e) $R$ is anti-symmetric since it only contains one element,
$(\mbox{Barack Obama}, \mbox{George W. Bush})$, and
(George W. Bush, Barack Obama)$\not\in R$.
\end{Answer}
\begin{Answer}{3.173}
(a) No.  The relation $R=\{(1,2),(2,1),(1,3)\}$ is neither symmetric ($(3,1)\not\in R$)
nor anti-symmetric ($(1,2)$ and $(2,1)$ are both in $R$).
(b) No.  For example, $R$ from answer (a) is not anti-symmetric, but isn't symmetric either.
(c) Yes.  If you answered incorrectly, don't worry.  You get to think about why the answer
is `yes' in the next exercise.
\end{Answer}
\begin{Answer}{3.174}
Many answers will work, but they all have the same thing in common: They only contain
`diagonal' elements (but not necessarily all of the diagonal elements).
For instance, let $R=\{(a,a): a\in \BBZ\}$.
Go back to the definitions for symmetric and anti-symmetric and verify that this is indeed both.
Another examples is $R=\{(\mbox{Ken}, \mbox{Ken})\}$ on the set of English words.
\end{Answer}
\begin{Answer}{3.176}
(a) $T$ is transitive since if $a$ is taller than $b$, and $b$ is taller than $c$, clearly $a$ is taller
than $c$.  In other words $(a,b)\in R$ and $(b,c)\in R$ implies that $(a,c)\in R$.
(b) $N$ is transitive because if $a$'s name starts with the same letter as $b$'s name, and
$b$'s name starts with the same letter as $c$'s name, clearly it is the same letter in all of them,
so $a$'s name starts with the same letter as $c$'s.
(c) $C$ is {\bf not} transitive.  You might think a similar argument as in (a) and (b) works here, but
it doesn't.  The proof from (b) works because names start with a single letter, so
transitivity holds.  But if $(a,b)\in C$ and $(b,c)\in C$, it might be because $a$ and $b$ have both
been to Chicago, and $b$ and $c$ have both been to New York, but that $a$ has never been to New York.
In this case, $(a,c)\not\in C$.  So $C$ is not transitive.
(d) $K$ is not transitive.  For instance, $(\mbox{David Letterman},\mbox{Chuck Cusack})\in K$
and $(\mbox{Chuck Cusack},\mbox{David Letterman's son})\in K$, but $(\mbox{David Letterman},\mbox{David Letterman's son})\not\in K$ since I sure hope he knows his own son.
(e) $R$ is transitive since there isn't even an $a,b,c\in R$ such that $(a,b)$ and $(b,c)$ are both in $R$,
so it holds vacuously.
\end{Answer}
\begin{Answer}{3.179}
(a) $T$ is {\bf not} an equivalence relation since it is not symmetric.
(b) $N$ is an equivalence relation since it is reflexive, symmetric, and transitive.
(c) $C$ is {\bf not} an equivalence relation since it is not transitive.
(d) $K$ is {\bf not} an equivalence relation since it is not reflexive, symmetric, or transitive.
This one isn't even close!
(e) $R$ is {\bf not} an equivalence relation since it is not reflexive.
\end{Answer}
\begin{Answer}{3.181}
(a) $T$ is a {\bf not} partial order because it is not reflexive.
(b) $N$ is {\bf not} a partial order since it is not anti-symmetric.
(c) $C$ is {\bf not} a partial order since it is not anti-symmetric or transitive.
(d) $K$ is {\bf not} a partial order since it is not reflexive, anti-symmetric, or transitive.
(e) $R$ is {\bf not} a partial order since it is not reflexive.
\end{Answer}
\begin{Answer}{3.182}
In the following, $A$, $B$, and $C$ are elements of $X$.  As such, they are sets.

\noindent
({\bf Reflexive})
Since $A\subseteq A$, $(A,A)\in R$, so $R$ is reflexive.

\noindent
({\bf Anti-symmetric})
If $(A,B)\in R$ and $(B,A)\in R$, then we know that
$A\subseteq B$ and $B\subseteq A$.  By Theorem~\ref{sets:containtmentproof}, this implies
that $A=B$.  Therefore $R$ is anti-symmetric.

\noindent
({\bf Transitive})
If $(A,B)\in R$ and $(B,C)\in R$, then $A\subseteq B$ and $B\subseteq C$.  But the definition
of $\subseteq$ implies that $A\subseteq C$, so $(A,C)\in R$, and $R$ is transitive.

\noindent
Since $R$ is reflexive, anti-symmetric, and transitive, it is a partial order.
\end{Answer}
\begin{Answer}{3.183}
(a) Since $(1,1)\not\in R$, $R$ is not reflexive.
(b) Since $(1,2)\in R$, but $(2,1)\not\in R$, $R$ is not symmetric.
(c) A careful examination of the elements reveals that it {\em is} anti-symmetric.
(d) A careful examination of the elements reveals that it {\em is} transitive.
(e) Since it is not reflexive or symmetric, it is not an equivalence relation.
(f) Since it is not reflexive, it is not a partial order.
\end{Answer}
\begin{Answer}{3.185}
$( (a,b), (a,b))$;
$bc$;
$da$;
$((c,d),(a,b))$;
symmetric;
$ad=bc$;
$cf=de$;
$de/f$;
$b(de/f)$;
$af=be$;
$( (a,b),(e,f) )$
\end{Answer}
